% ===================================================================
%  TABULKA č. 4 — Zralý věk a stáří.
%  Traduction 2026 d'apres texto/io, controlee sur texto/fr, texto/ru
%  et texto/uk. Consigne : voir texto/cs/10-tabulka-01.tex.
%
%  LA PARENTE PAR ALLIANCE EST LE LIEU OU LE TCHEQUE EST LE PLUS
%  ANCIEN DE TOUT LE VOLUME. La ou les langues romanes prefixent un
%  element au nom — « beau-pere », « suegro », « sogro » — et ou
%  l'anglais dit « -in-law », le tcheque a des mots SIMPLES, herites
%  sans composition : « tchán » le beau-pere, « tchyně » la
%  belle-mere, « zeť » le gendre, « snacha » la bru, « švagr » le
%  beau-frere, « švagrová » la belle-soeur. Six degres, six mots, pas
%  un compose.
%
%  ET LES MALADIES DE LA SECONDE SCENE FONT L'INVERSE, comme en
%  galicien : « revmatismus », « zápal plic », « horečka »,
%  « mdloby », « lék ». Le corps et la famille s'apprennent a la
%  maison ; la maladie s'apprend chez le medecin, et le medecin parle
%  latin en Boheme comme en Galice.
% ===================================================================

%%K t04-apar-1 apar
\VUsaut{4.0mm}
\VUtitre{15.0pt}{60}{TABULKA č. 4}
\VUsaut{1.6mm}
\VUtitre{11.4pt}{0}{Zralý věk a stáří.}
\VUsaut{3.2mm}

%%K t04-tit-1 sub
\VUcentre{11.4pt}{0}{\textit{První scéna.}}
\VUsaut{1.6mm}
\VUtitre{11.4pt}{0}{Svatební hostina.}
\VUsaut{2.6mm}

%%K t04-tit-2 sub
\VUpk{10.2pt}{40}{Podzim.}
\VUsaut{3.0mm}

%%K t04-c1-01-1 p
1. --- Mladí lidé vyrostli. Jeden z nich, Jan, se žení. Jsme na \VUgras{svatební hostině},
která shromažďuje kolem téhož stolu rodiny novomanželů.

%%K t04-c1-02-1 p
2. --- Jsme na \VUgras{podzim,} v měsíci \VUgras{září.} Studený vítr \VUgras{října} a
\VUgras{listopadu} ještě neobral venkov o jeho zelené listí. Večerní vzduch je vlažný
a provoněný; proto se večeře koná pod \VUgras{verandou} \textsuperscript{(8)} u zahrady.

%%K t04-c1-03-1 p
3. --- Daleko, na \VUgras{pahorku} \textsuperscript{(3)}, stojí domov Janových rodičů,
elegantní \VUgras{zámek} \textsuperscript{(1)} ukrytý v zeleni; krásné vinice, které dávají
výborné víno, pokrývají svah pahorku. Vinař je pěstuje a pečuje o ně.

%%K t04-c1-04-1 p
4. --- Nad vinicemi přelétá hejno \VUgras{koroptví} \textsuperscript{(2)} vyplašené příchodem
\VUgras{hajného} \textsuperscript{(5)}. Ten, s \VUgras{puškou} pod paží a s
\VUgras{brašnou na zvěř přes rameno}, prochází \VUgras{houštinu}
\textsuperscript{(4)} se svým \VUgras{psem} \textsuperscript{(6)}. Dohlíží na panství a brání
pytlákům hubit zvěř.

%%K t04-c1-05-1 p
5. --- Vně \VUgras{verandy} šplhá \VUgras{vinná réva,} z níž visí husté \VUgras{hrozny}
\textsuperscript{(7)}.

%%K t04-c1-06-1 p
6. --- Krásné podzimní květiny, které zdobí \VUgras{stůl} \textsuperscript{(16)}, jsou
\VUgras{chryzantémy} \textsuperscript{(9)}; \VUgras{otec} \textsuperscript{(10)}
nevěsty si jednu z nich vložil do knoflíkové dírky svého fraku.

%%K t04-c1-07-1 p
7. --- Jídlo se chýlí ke konci. \VUgras{Sluha} \textsuperscript{(11)} začíná přinášet zákusky
a brzy odnese různé části příboru, \VUgras{lžíce} \textsuperscript{(14)},
\VUgras{vidličky} \textsuperscript{(12)}, \VUgras{nože} \textsuperscript{(13)},
\VUgras{mísy} \textsuperscript{(15)}, jichž už není třeba.

%%K t04-tit-3 sub
\VUpk{10.2pt}{40}{Příbuzenstvo.}
\VUsaut{3.0mm}

%%K t04-c1-08-1 p
8. --- Všichni vypadají vesele, a úsměv, s nímž \VUgras{matka} \textsuperscript{(17)} ženicha
hledí na své děti, dokazuje její štěstí.

%%K t04-c1-09-1 p
9. --- Ta svatba spojuje dvě bohaté rodiny kraje. Jan, \VUgras{manžel} \textsuperscript{(18)},
je \VUgras{synem} velkého statkáře a stává se \VUgras{zetěm} obratného továrníka.

%%K t04-c1-10-1 p
10. --- \VUgras{Manželé} jsou mladí, a přesto byli dlouho \VUgras{zasnoubeni.}
\VUgras{Mladá žena} \textsuperscript{(19)}, Marta, je půvabná ve svých \VUgras{bílých
šatech} ozdobených pomerančovými květy.

%%K t04-c1-11-1 p
11. --- Její \VUgras{tchán} \textsuperscript{(20)} je velmi šťasten, že má tak laskavou
\VUgras{snachu,} a Janova \VUgras{tchyně} \textsuperscript{(21)} se usmívá, když vidí svou
\VUgras{dceru} tak krásnou.

%%K t04-c1-12-1 p
12. --- Na konci stolu sedí ostatní \VUgras{hosté} \textsuperscript{(22)},
\VUgras{příbuzní, přátelé, bratranci, sestřenice.} \VUgras{Vrchní}
\textsuperscript{(23)} se servisem pod paží je pilně obsluhuje.

%%K t04-tit-4 sub
\VUpk{10.2pt}{40}{Přátelé.}
\VUsaut{3.0mm}

%%K t04-c1-13-1 p
13. --- Napravo od stolu, zde je \VUgras{slečna} \textsuperscript{(24)},
\VUgras{přítelkyně} mladé nevěsty, potom její \VUgras{bratr}, který je jejím
\VUgras{družbou} \textsuperscript{(25)}. Rozmlouvá se svou \VUgras{družičkou}
\textsuperscript{(26)}, ženichovou sestrou, která se touto svatbou stává Martinou
\VUgras{švagrovou.} Je to půvabná dívka. Má krásné \VUgras{šperky, perlový náhrdelník}
a zlatý \VUgras{náramek.}

%%K t04-c1-14-1 p
14. --- Blízko nich, pán, který stojí a zvedá svou sklenici, je \VUgras{přítelem}
\textsuperscript{(27)} rodiny. Je jedním ze \VUgras{ženichových svědků.}
\VUgras{Připíjí,} pije na \VUgras{zdraví} novomanželů a přeje jim dlouhé štěstí. Je
oblečen ve slavnostním obleku, ve fraku s \VUgras{lakovými botami} \textsuperscript{(28)}. Je
průvodcem \VUgras{babičky} \textsuperscript{(29)}, která je
\VUgras{vdova.}

%%K t04-c1-15-1 p
15. --- Mladík ve \VUgras{fraku} \textsuperscript{(31)}, s tenkým černým knírkem, je Janův
\VUgras{švagr} \textsuperscript{(30)}. Je vynikající inženýr. Je sousedem velmi elegantní
\VUgras{mladé dámy} \textsuperscript{(33)}, Martiny \VUgras{starší sestry} a matky roztomilé
holčičky, která přichází, podávajíc ruku svému bratranci, nabídnout květinu a \VUgras{popřát}
jí \VUgras{dobrý večer.} Ta dáma má velmi krásné
\VUgras{náušnice} a elegantní \VUgras{čepec.}

%%K t04-c1-16-1 p
16. --- Malý bratranec, tak podivný ve své \VUgras{halence} \textsuperscript{(36)} a se svým
nadouvajícím se \VUgras{rukávem} \textsuperscript{(37)}, je velmi politováníhodný. Je
\VUgras{sirotek} \textsuperscript{(35)}. Velmi mladý ztratil otce i matku. Jeho strýc je jeho
\VUgras{poručníkem} \textsuperscript{(38)} a zůstal svobodný, aby ubohé dítě vychoval. Je to
dobrotivý muž. Je poněkud plešatý a velmi krátkozraký; užívá
\VUgras{skřipec.}

%%K t04-tit-5 sub
\VUsaut{2.4mm}
\VUpk{10.2pt}{40}{Zákusky.}
\VUsaut{3.0mm}

%%K t04-c1-17-1 p
17. --- Za hodovníky, na stole pokrytém \VUgras{ubrusem,} jsou připraveny
\VUgras{zákusky} \textsuperscript{(39)}. Ve velké míse, zde je ovoce,
\VUgras{broskve} \textsuperscript{(40)}, \VUgras{hrušky} \textsuperscript{(41)},
\VUgras{jablka} \textsuperscript{(42)}, \VUgras{meruňky} \textsuperscript{(43)} a
\VUgras{hrozny} \textsuperscript{(44)}. Blízko, \VUgras{ananasy}
\textsuperscript{(45)}, krásný \VUgras{dort} \textsuperscript{(46)}, \VUgras{láhve likéru}
\textsuperscript{(48)} a \VUgras{šampaňské} \textsuperscript{(49)}, a také sloupce čistých
\VUgras{talířů} \textsuperscript{(47)}.

%%K t04-c1-18-1 p
18. --- \VUgras{Sklepník} \textsuperscript{(50)} odzátkovává \VUgras{láhev}
\textsuperscript{(52)} starého vína \VUgras{vývrtkou} \textsuperscript{(51)}.
\VUgras{Korková zátka} \textsuperscript{(53)} se zdá velmi těžko vytažitelná. Ten
sklepník má \VUgras{servítek} \textsuperscript{(54)} pod paží.

%%K t04-c1-19-1 p
19. --- Po večeři půjdou pánové do kuřárny a dámy se půjdou připravit na ples.

%%K t04-tit-6 sub
\VUsaut{5.0mm}
\VUcentre{11.4pt}{0}{\textit{Druhá scéna.}}
\VUsaut{1.6mm}
\VUpk{11.4pt}{40}{Dědečkovy jmeniny.}
\VUsaut{2.6mm}

%%K t04-tit-7 sub
\VUpk{10.2pt}{40}{Návštěva.}
\VUsaut{3.0mm}

%%K t04-c2-01-1 p
1. --- Uplynulo deset let; Janovi rodiče zestárli a mají velmi rádi své tři vnoučata, dva
\VUgras{vnuky} \textsuperscript{(2)} a jednu \VUgras{vnučku} \textsuperscript{(3)}. Protože
jsou dnes \VUgras{dědečkovy jmeniny,} přicházejí děti popřát jim hezký svátek.

%%K t04-c2-02-1 p
2. --- Malý Petr je ještě velmi malý; jsou mu teprve tři roky. Vidí, jak se blíží
\VUgras{kočka} \textsuperscript{(1)}. Chce pohladit její krásnou hedvábnou
\VUgras{srst} a její dlouhý \VUgras{ocas}; nepomyslí na to, že se pod půvabnými
\VUgras{tlapkami} skrývají zlomyslné ostré drápy, které ho snad poškrábou.

%%K t04-c2-03-1 p
3. --- Jeho sestra Gabriela, která mu podává ruku, je mnohem vážnější, a Marcel se svým velkým
\VUgras{pláštěm} \textsuperscript{(4)} a koženým \VUgras{opaskem} \textsuperscript{(5)} vypadá
poněkud zmužile, přes svůj krátký rukáv, který nechává vidět jeho \VUgras{punčochy}
\textsuperscript{(6)}.

%%K t04-c2-04-1 p
4. --- Jejich otec si oblékl svůj tlustý zimní \VUgras{svrchník} \textsuperscript{(7)} s
\VUgras{kožešinovým límcem} \textsuperscript{(8)} a ve své
\VUgras{kapse} \textsuperscript{(9)} má šátek. Jeho žena má svůj plstěný klobouk
ozdobený \VUgras{péry} \textsuperscript{(10)}, svůj \VUgras{kabátek} \textsuperscript{(11)},
svou \VUgras{kožešinovou boa} \textsuperscript{(12)} a svůj
\VUgras{rukávník} \textsuperscript{(13)}.

%%K t04-c2-05-1 p
5. --- Je velká zima, protože přišel mráz. \VUgras{Sníh} \textsuperscript{(19)} padal celý den
a střechy domů jsou zcela bílé. Zároveň s \VUgras{nocí} se chlad stává ještě ostřejším. Na
nebi svítí \VUgras{měsíc} \textsuperscript{(17)} a \VUgras{hvězdy} \textsuperscript{(18)} a
pronikají \VUgras{temnotou.}

%%K t04-tit-8 sub
\VUsaut{4.2mm}
\VUpk{10.2pt}{40}{Nemoc.}
\VUsaut{3.0mm}

%%K t04-c2-06-1 p
6. --- Ubohý \VUgras{slepec} \textsuperscript{(20)}, který přechází náměstí před
\VUgras{lékárnou} \textsuperscript{(16)} veden svým \VUgras{pudlem}
\textsuperscript{(21)}, je velmi nešťastný. Justýna, \VUgras{služka} \textsuperscript{(14)},
přináší z kuchyně velký \VUgras{hrnek} \textsuperscript{(15)} velmi horkého
\VUgras{bylinného čaje} štědře oslazeného.

%%K t04-c2-07-1 p
7. --- \VUgras{Babička} \textsuperscript{(23)}, která chce hledat na stole svůj
\VUgras{lorňon} \textsuperscript{(26)}, aby přečetla \VUgras{předpis}
\textsuperscript{(27)}, který \VUgras{lékař} \textsuperscript{(25)} právě napsal, vstává ze
svého křesla opírajíc se o svou \VUgras{hůl} \textsuperscript{(24)}.

%%K t04-c2-08-1 p
8. --- Často trpí \VUgras{neduhy, závratěmi, nachlazeními, bolestmi zubů}; její nohy jsou
slabé a její oči už nejsou příliš dobré. Přesto si ráda tropí žerty ze svého starého
\VUgras{lékaře,} který jí stále chce předepisovat nějaký \VUgras{lék} \textsuperscript{(28)}.
Dnes je velmi veselá, protože má kolem sebe svou mladou rodinu.

%%K t04-c2-09-1 p
9. --- V dobře zavřeném pokoji je pohodlně. V mramorovém \VUgras{krbu} \textsuperscript{(29)}
plápolá \VUgras{překrásný oheň} \textsuperscript{(30)} a člověk se cítí šťastný v mírném
\VUgras{světle} \VUgras{lampy} \textsuperscript{(31)}, kterou právě rozsvítili.

%%K t04-tit-9 sub
\VUsaut{4.2mm}
\VUpk{10.2pt}{40}{Dědeček.}
\VUsaut{3.0mm}

%%K t04-c2-10-1 p
10. --- I \VUgras{dědeček} \textsuperscript{(33)} zapomněl, že byl nemocný, že ho jeho
\VUgras{revmatismus} připoutává k jeho \VUgras{křeslu} \textsuperscript{(34)} a že
ještě včera měl \VUgras{horečku a mdloby.} Už si nevzpomíná, že loni v zimě trpěl
\VUgras{zápalem plic} a že chraptivý a mučivý \VUgras{kašel} ho velmi soužil.

%%K t04-c2-10-2 p
A přece se uzdravil jen velmi obtížně. Od té doby je mu poněkud lépe. Ale noha, kterou si
opírá o \VUgras{polštář} \textsuperscript{(38)} položený na malé \VUgras{stoličce}
\textsuperscript{(39)}, ho také bolí.

%%K t04-c2-11-1 p
11. --- Když je pečlivě zabalen do svého teplého \VUgras{županu} \textsuperscript{(35)} s
tlustými perleťovými \VUgras{knoflíky} \textsuperscript{(36)} a když má blízko sebe, aby mu
četla noviny, malou \VUgras{sirotu} \textsuperscript{(40)} oblečenou v bílém \VUgras{čepci,} v
modrých \VUgras{šatech} \textsuperscript{(42)} a v \VUgras{pelerínce} \textsuperscript{(41)},
nestěžuje si.

%%K t04-c2-12-1 p
12. --- Každý rok, když \VUgras{kalendář} \textsuperscript{(43)} znovu ukazuje
\VUgras{datum} prvního \VUgras{prosince,} netrpělivě očekává svá \VUgras{vnoučata,}
protože ví, že na to důležité \VUgras{datum} nezapomenou.

%%K t04-c2-13-1 p
13. --- S dojetím vzpomíná na dávnou dobu, kdy sám přišel popřát hezký svátek slavnému
císařskému \VUgras{maršálovi,} jehož \VUgras{podobizna} \textsuperscript{(44)} visí na stěně a
který je \VUgras{pradědem} jeho dětí.
